% Part I, chapter 1: Introduction.
\chapter{Introduction}
\label{chap:introduction}

\wfdbx{} is a reimplementation of the WFDB Software Package, version
\wfdbversion{}: the C library and the command-line tools with which
physiologic signal and annotation files in the WFDB format, such as those of
the PhysioNet databases, are read, written and converted. \wfdbx{}
reimplements the package in seven programming languages: Python, C++, Rust,
Julia, R, MATLAB/Octave and Kotlin. In each of them it reads and writes the
same files as the original C library, provides the same functions under
their C names, and provides the same command-line tools, with the same
options, output, messages and exit status. The seven implementations are
tested to give identical results, byte for byte, to the original C programs
and to each other.

\wfdbx{} is an independent project. It is not part of the WFDB project
(\url{https://wfdb.io/}) or of PhysioNet (\url{https://physionet.org/}), and
is not endorsed by them. It follows the WFDB format and the behavior of the
original C package, which serves as its specification: the expected results
of all its tests were produced by the C programs.

\section{What WFDB-x is}
\label{sec:what}

Each of the seven implementations
\begin{itemize}
\item reads and writes exactly the same files as the original C library:
  headers (\file{.hea}), signal files in every WFDB storage format (8, 16,
  61, 80, 160, 212, 310, 311, 24, 32, and the FLAC formats 508, 516 and
  524), annotation files (in MIT and AHA format), EDF and EDF+ files,
  calibration files, info files, and multi-segment records (see
  \cref{chap:records});
\item provides the same 33 command-line tools as the C package (listed in
  \cref{sec:three-ways}), with the same options;
\item provides the WFDB library functions under their C names
  (\code{isigopen}, \code{getvec}, \code{getann}, \code{putann},
  \code{strtim}, \ldots), with the behavior of the C functions, and a
  high-level API in the idiom of the language;
\item is written natively in its language: none of the implementations calls
  the C library or another language. It needs nothing beyond the standard
  library of its language, except NumPy for Python. FLAC files are read and
  written without libFLAC. Records are read from URLs with the system
  \cmd{curl} program in C++ and Rust (and in R and MATLAB/Octave when
  \cmd{curl} is installed), and with the standard facilities of the
  language otherwise (see \cref{sec:network}).
\end{itemize}

\wfdbx{} behaves exactly as the original package, with a few deliberate
exceptions, which are the same in all seven languages and are tested like
everything else: it never assumes a sampling frequency, gain, ADC resolution
or units that a record does not give; it reads and writes files only where
they are named, and opens URLs only when network access is turned on; it
sorts annotation files itself, without losing annotations; a tool never
writes a signal file over a signal file it is reading; and a tool that cannot
write its output exits with status 1. \Cref{chap:behavior} describes each of
them, and \emph{legacy mode}, which restores the original behavior exactly.

The interactive viewer of the C package (WAVE), its plotting programs, and
the analysis programs of its \file{psd/} directory are outside the scope of
\wfdbx{}.

The speed and memory use of the seven implementations, compared with those
of the C library and tools, are measured in \file{docs/BENCHMARKS.md}. In
short, Rust, Julia and C++ read an hour of a two-signal ECG in 4 to 6\,ms,
4 to 7 times faster than the C library, which reads one frame per
\code{getvec} call; the high-level functions of every language read whole
blocks of samples, and are the fastest way to read a record in that
language.

\section{The seven implementations}
\label{sec:implementations}

All seven implement the complete library and all 33 tools. Each has a
directory of the repository, with a \file{README.md} that gives its
requirements, installation instructions and examples, and a chapter of
\cref{part:languages}. \Cref{tab:implementations} shows how each is
installed, from a copy of the repository (or from PyPI for Python), and how
a tool is run in it. In code, the packages are named \code{wfdbx} (Python,
MATLAB/Octave, Rust, R, Kotlin, and the C++ namespace) or \code{WFDBX}
(Julia).

\begin{table}[htbp]
\caption{The seven implementations: installation, and how a tool is run.}
\label{tab:implementations}
\begin{tabularx}{\linewidth}{@{}P{2.3cm}P{1.5cm}LP{3.0cm}@{}}
\toprule
{\tabhead Language} & {\tabhead Directory} & {\tabhead Installation} &
  {\tabhead Running a tool} \\
\midrule
Python & \file{python/} & \cmd{pip install wfdb-x} (PyPI), or
  \cmd{pip install ./python} & \cmd{wfdbx rdsamp -r 100s} \\
C++ & \file{cpp/} & \cmd{cmake -S cpp -B cpp/build \&\& cmake --build
  cpp/build} & \cmd{rdsamp -r 100s} \\
Rust & \file{rust/} & \cmd{cargo install --path rust} &
  \cmd{rdsamp -r 100s} \\
Julia & \file{julia/} & \cmd{using Pkg;
  Pkg.develop(path="/path/to/wfdb-x/julia")} & \cmd{wfdbx rdsamp -r 100s} \\
R & \file{r/} & \cmd{R CMD INSTALL r} & \cmd{wfdbx rdsamp -r 100s} \\
MATLAB / GNU Octave & \file{matlab/} &
  \cmd{addpath('/path/to/wfdb-x/matlab')} (tested with GNU Octave; MATLAB
  untested) & \cmd{wfdbx rdsamp -r 100s} \\
Kotlin (Android, iOS, JVM) & \file{kotlin/} & \cmd{cd kotlin \&\&
  ./gradlew publishToMavenLocal}, then
  \cmd{implementation("io.github.\allowbreak alphanumericslab:\allowbreak wfdbx:1.2.1")} (Gradle,
  with \code{mavenLocal()}) & \cmd{kotlin/bin/wfdbx rdsamp -r 100s} \\
\bottomrule
\end{tabularx}
\end{table}

In Julia, R and MATLAB/Octave, \cmd{wfdbx} is a short alias for the
language's own launcher of the tools; the chapter of each language in
\cref{part:languages} shows how to set it up. The Python package is
distributed on PyPI as \code{wfdb-x} and installs the \cmd{wfdbx} command;
\cmd{wfdbx --list} lists the tools, \cmd{wfdbx --version} prints the version,
and \cmd{wfdbx --examples} copies the example notebooks of the package.

The Kotlin implementation runs on the JVM, on Android, and on iOS through
Kotlin Multiplatform; it is tested on the JVM. The MATLAB/Octave
implementation is written for both MATLAB and GNU Octave, and is tested with
GNU Octave.

\paragraph{The sample record.} Record \code{100s}, the first minute of record
100 of the MIT-BIH Arrhythmia Database (two ECG signals sampled at 360\,Hz,
and its beat annotations), is in the directory \file{tests/data} of the
repository, and the Python package includes a copy of it (in the directory
that \code{wfdbx.sample\_data\_dir()} returns). The examples of this manual
use it. To run them, run the tools in that directory, or set the input
directory to it (see \cref{chap:files}), for example:
\begin{lstlisting}[style=shell]
cd /path/to/wfdb-x/tests/data
wfdbx rdsamp -r 100s -t s3
\end{lstlisting}
which prints the first three samples of each signal, in ADC units:
\begin{lstlisting}[style=text]
              0	    995	   1011
              1	    995	   1011
              2	    995	   1011
\end{lstlisting}

\section{Three ways to use WFDB-x}
\label{sec:three-ways}

Every language offers the same three ways in, summarized in
\cref{tab:three-ways}. All three are written natively in the language.

\begin{table}[htbp]
\caption{The three ways to use \wfdbx{}.}
\label{tab:three-ways}
\begin{tabularx}{\linewidth}{@{}P{2.8cm}LL@{}}
\toprule
& {\tabhead What it is} & {\tabhead Use it for} \\
\midrule
1.~The high-level API (recommended) &
  A few functions in the idiom of the language: \code{read\_record},
  \code{read\_header}, \code{read\_annotations}, \code{write\_record},
  \code{write\_annotations}, \code{check\_record}, and the header metadata
  functions \code{read\_metadata}, \code{get\_metadata},
  \code{set\_metadata}, \code{remove\_metadata} and \code{update\_metadata}
  (\code{readRecord}, \ldots{} in Kotlin). One call reads a whole record or
  annotation file into arrays, with named options, errors reported as
  exceptions or error values, and fast block reads. It has no counterpart in
  the C package. &
  New code, analysis, notebooks, machine learning: anything that needs a
  record as arrays. \\
\addlinespace
2.~The WFDB library functions &
  The C library API under its C names (\code{isigopen}, \code{getvec},
  \code{getann}, \code{putann}, \code{strtim}, \ldots), with the C behavior,
  frame by frame and annotation by annotation. &
  Porting a program written for the C library, or processing a record as a
  stream with the exact C semantics. \\
\addlinespace
3.~The command-line tools &
  The 33 tools (\cmd{rdsamp}, \cmd{rdann}, \cmd{wrsamp}, \cmd{xform},
  \ldots) with the options, output, messages and exit status of the C
  tools, implemented in each language. &
  Shell scripts and pipelines, and reproducing work done with the original
  tools, without the C package. \\
\bottomrule
\end{tabularx}
\end{table}

The high-level API is built on the library functions, and returns exactly
the samples and annotations that they, and the C library, give; the tools are
built on the library functions too. The three can be mixed in one program,
with one caution: the high-level readers and writers call \code{wfdbquit},
which closes the files opened with the library functions.

In Python, for example, the first three samples of record \code{100s} are
read in physical units with the high-level API as follows:
\begin{lstlisting}[language=Python,style=example]
import wfdbx
rec = wfdbx.read_record("100s", sampto=3, physical=True)
print(rec.fs, rec.names, rec.signals.shape)   # 360.0 ['MLII', 'V5'] (3, 2)
print(rec.signals[0])                         # [-0.145 -0.065]
\end{lstlisting}
and the first sample of each signal is read in ADC units, and converted, with
the library functions:
\begin{lstlisting}[language=Python,style=example]
import wfdbx
si = []
n = wfdbx.isigopen("100s", si, 2)     # 2 signals, described in si
v = [0] * n
wfdbx.getvec(v)                       # v is [995, 1011]
print(wfdbx.aduphys(0, v[0]))         # -0.145
wfdbx.wfdbquit()
\end{lstlisting}

The 33 tools are those of the \file{app/} and \file{convert/} directories of
the C package:
\begin{itemize}
\item signals and annotations: \cmd{rdsamp}, \cmd{wrsamp}, \cmd{rdann},
  \cmd{wrann}, \cmd{sortann}, \cmd{ann2rr}, \cmd{rr2ann}, \cmd{wfdbdesc},
  \cmd{wfdbtime}, \cmd{time2sec}, \cmd{sampfreq}, \cmd{signame},
  \cmd{signum}, \cmd{sigamp}, \cmd{snip}, \cmd{xform}, \cmd{wfdbcat},
  \cmd{wfdbwhich};
\item format converters: \cmd{edf2mit}, \cmd{mit2edf}, \cmd{rdedfann},
  \cmd{wav2mit}, \cmd{mit2wav}, \cmd{wfdb2mat}, \cmd{a2m}, \cmd{m2a},
  \cmd{ad2m}, \cmd{md2a}, \cmd{ahaecg2mit}, \cmd{makeid}, \cmd{readid},
  \cmd{revise}, \cmd{parsescp}.
\end{itemize}
\Cref{part:reference} describes each of them, and each library function.

\section{How the results are checked}
\label{sec:checked}

For the same input, a \wfdbx{} tool produces byte for byte the same standard
output, standard error (including its error messages), exit status and
output files as the original C tool, and the library functions return the
same values as their C counterparts. The exceptions are the deliberate
differences of \cref{chap:behavior}; the reference results are produced by
the C code with these changes applied, so that they are tested in the same
way. This is verified by tests whose expected results were produced by the
original C code, not written by hand:
\begin{itemize}
\item \textbf{The conformance suite} (\file{tests/conformance/}, described in
  \file{docs/conformance.md}): 1135 command-line test cases, run by a
  harness that does not depend on the language of the implementation. Their
  expected results were produced by the C programs, built from the original
  sources (\file{third\_party/wfdb-main.zip}) with FLAC support and with
  the patch \file{tests/reference/divergences.patch}, which defines the
  deliberate differences, applied. The exceptions are 20 cases whose input
  makes the C programs crash; their expected results are the behavior that
  \wfdbx{} defines for that input (\cref{sec:crashes}). A second suite, for
  legacy mode, has 883 cases: the cases of the original package, run with
  \opt{--legacy} and compared with the results of the unmodified C package.
  Every implementation passes all the cases of both suites on the reference
  platform, x86-64 Linux with glibc.
\item \textbf{Test records} (\file{tests/data/}), written by the C library
  itself (\file{tests/reference/gen\_records.c}). They cover every storage
  format, multi-frequency records, multi-segment records of fixed and of
  variable layout, skewed signals, byte offsets, block-structured files, EDF,
  and annotation files that use every feature of the annotation format.
\item \textbf{Unit tests} of the library and of the high-level functions, in
  each language directory, with expected values printed by the C library
  (\file{tests/reference/api\_expect.c}).
\item \textbf{Metadata tests} (\file{tests/metadata/}): 49 cases of the
  header metadata functions, which every implementation must pass with
  byte-identical headers, and a check that the original C programs read the
  changed headers as before (\cref{chap:metadata}).
\end{itemize}
Because all seven implementations must pass the same cases, each is also
checked against the other six. The MATLAB/Octave code is tested with GNU
Octave, also through the MATLAB driver \file{matlab/bin/wfdbx\_matlab.sh};
MATLAB itself has not been available to run the tests.

A few behaviors of the C code are undefined (crashes, reads of uninitialized
memory), and a few depend on the platform. They are listed, with the
behavior that \wfdbx{} chose, in \file{docs/KNOWN\_C\_BEHAVIORS.md} and
summarized in \cref{sec:known-c}.

\section{Versions}
\label{sec:versions}

\wfdbx{} has its own version numbers, shared by the seven implementations,
which follow semantic versioning (\url{https://semver.org}): a new patch
number (1.1.x) for corrections only, a new minor number (1.x.0) for new
features that do not change existing behavior, and a new major number
(x.0.0) for changes that can break existing code. This manual describes
version \wfdbxversion{}. The file \file{CHANGELOG.md} lists the changes in
each version:
\begin{itemize}
\item 1.0.0 (16 October 2026): the WFDB library, its 33 command-line tools
  and a high-level API in Python, C++, MATLAB/Octave, Rust, Julia and R;
\item 1.0.1 (27 October 2026): the Kotlin Multiplatform implementation, the
  seventh language;
\item 1.1.0 (28 October 2026): faster reading and writing of records, signal
  names without surrounding spaces in the high-level readers, and write
  errors and frame checks in the high-level writers;
\item 1.2.0: the header metadata functions (\cref{chap:metadata}), this
  manual, and the reference to cite (\cref{sec:cite}). No existing function
  or tool changes behavior, except for three corrections: the files that a
  tool reads or creates itself follow the input and output directories,
  \code{osigopen(record, NULL, 0)} returns 0 in every language, and the
  high-level functions leave the annotation time resolution and the getvec
  mode as they were (\cref{sec:corrections-1-2-0});
\item 1.2.1: the library under the GNU Lesser General Public License,
  version 2.1 or later, with the command-line tools still under the GNU
  General Public License, version 3 or later (\cref{sec:license}); and the
  C++ library builds on Windows again.
\end{itemize}

The version numbers of \wfdbx{} are independent of the version of the WFDB
Software Package that it implements, \wfdbversion{}, in every version. The
library reports that version as the C library does: \code{wfdbversion()}
returns \code{"10.7.0"}, and \code{wfdberror()}, before any error, returns
\code{"WFDB library version 10.7.0 (Jun 10 2022).\textbackslash{}n"}.

\section{How this manual is organized}
\label{sec:organization}

\begin{itemize}
\item \textbf{\Cref{part:concepts}, Concepts}, describes what the seven
  implementations share: this introduction; the WFDB records and their files
  (\cref{chap:records}); where files are read and written, network access
  and the environment variables (\cref{chap:files}); the differences from
  the original package, legacy mode, and the reporting of errors
  (\cref{chap:behavior}); and the header metadata (\cref{chap:metadata}).
\item \textbf{\Cref{part:languages}, The seven languages}, has one chapter
  for the users of each language: installation, the high-level API with
  every function, its arguments, results, errors and examples, and how the
  library functions and the tools are used in that language
  (\cref{chap:python,chap:cpp,chap:rust,chap:julia,chap:r,chap:matlab,chap:kotlin}).
\item \textbf{\Cref{part:reference}, Reference}, describes the WFDB library
  functions and the command-line tools, which are the same in every
  language, and the storage formats and annotation codes
  (\cref{chap:formats}).
\end{itemize}

\section{Typographic conventions}
\label{sec:conventions}

\begin{itemize}
\item \code{read\_record}, \code{WFDB\_INVALID\_SAMPLE}, \code{"0:10"}:
  identifiers, values and code, in a typewriter font.
\item \cmd{rdsamp -r 100s}: a command line; \opt{-r}: an option of a tool;
  \file{tests/data/100s.hea}: a file or a directory.
\item The WFDB library functions and the tools are named by their C names,
  which are their names in every language.
\item Listings of code and of command lines are set on a shaded background.
  A listing on a white background, with a bold blue rule, is the call syntax
  of a function, copied from the code. The Python examples marked as
  examples are run, as part of the checks of this manual, in a directory
  that holds a copy of \file{tests/data}, which is the current directory and
  the input directory; their results are given in comments.
\end{itemize}
Three kinds of boxes set off remarks:
\begin{note}
A note gives details that matter in some cases, such as a behavior of the C
package that \wfdbx{} reproduces.
\end{note}
\begin{tip}
A tip gives a way to do something faster or more simply.
\end{tip}
\begin{phibox}
A box of this kind concerns data that can identify patients, such as names,
dates or ages stored in a header.
\end{phibox}

\section{How to cite}
\label{sec:cite}

To refer to \wfdbx{} in a publication, cite the preprint that describes it:
\begin{quote}
Reza Sameni. \emph{WFDB-x: Conformance-Tested Implementations of the WFDB
Software Package in Seven Programming Languages.} arXiv:2611.01358, 2026.\\
DOI: \href{https://doi.org/10.48550/arXiv.2611.01358}{10.48550/ARXIV.2611.01358}.\\
\url{https://arxiv.org/abs/2611.01358}
\end{quote}
As a BibTeX entry:
\begin{lstlisting}[style=text]
@misc{wfdb-x,
  doi = {10.48550/ARXIV.2611.01358},
  url = {https://arxiv.org/abs/2611.01358},
  author = {Sameni,  Reza},
  title = {{WFDB-x: Conformance-Tested Implementations of the WFDB Software Package in Seven Programming Languages}},
  publisher = {arXiv},
  year = {2026},
  copyright = {Creative Commons Attribution Share Alike 4.0 International}
}
\end{lstlisting}
The same reference is in the file \file{CITATION.cff} of the repository.
Work that uses the WFDB format or PhysioNet data should also cite the WFDB
Software Package and PhysioNet, as PhysioNet asks.

\section{Authors and license}
\label{sec:license}

The WFDB Software Package was written by George B.~Moody, Benjamin Moody and
the WFDB contributors; it is Copyright \textcopyright{} 1983--2022 George
B.~Moody and the WFDB contributors. WFDB was originally developed at MIT by
George B.~Moody, and builds on decades of work by him, Roger G.~Mark and the
broader research community. The authors of the C code that each part of
\wfdbx{} translates are credited, file by file, in each \wfdbx{} source file
and in \file{docs/AUTHORS.md}.

\wfdbx{} was developed by Reza Sameni (2026). It is an independent
reimplementation of the WFDB Software Package: its code in each language is a
translation of the package's C code, written for \wfdbx{}, not code by the
original WFDB developers.

The \wfdbx{} library is distributed under the GNU Lesser General Public
License, version 2.1 or (at the user's option) any later version
(LGPL-2.1-or-later; see the file \file{LICENSE}): programs under any license,
proprietary ones included, may use it on the LGPL's terms. The command-line
tools (each language's \code{apps} and the \cmd{wfdbx} command), translations
of the tools of the original package, are distributed under the GNU General
Public License, version 3 or (at the user's option) any later version
(GPL-3.0-or-later; see the file \file{LICENSE-GPL}). Each source file states
its license. The original C sources, kept unmodified in
\file{third\_party/}, keep their own licenses: the GNU Library General Public
License, version 2 or later, for the library, and the GNU General Public
License, version 2 or later, for the tools; their ``or later'' terms allow
those of \wfdbx{}.
